%% kapa: exam=1 round=34 year=2023 period=1 form=A subject=real_estate subject_ko=부동산학원론 seq=22 number=102
%% source: Q-Net 감정평가사 기출문제 (cst00309) artlSeq=5211789 file=raw/1008/5211789/2023년도 제34회 감정평가사 1차시험 1교시.pdf
%% conversion: kordoc 4.24.51 -> tools/gen1.py (auto); review=proofread-sonnet; flags=none
\begin{question}{102}
프로젝트 금융에 관한 설명으로 옳은 것은?
\begin{choices}
  \item 기업전체의 자산 또는 신용을 바탕으로 자금을 조달하고, 기업의 수익으로 원리금을 상환하거나 수익을 배당하는 방식의 자금조달기법이다.
  \item 프로젝트 사업주는 기업 또는 개인일 수 있으나, 법인은 될 수 없다.
  \item 프로젝트 사업주는 대출기관으로부터 상환청구를 받지는 않으나, 이러한 방식으로 조달한 부채는 사업주의 재무상태표에는 부채로 계상된다.
  \item 프로젝트 회사가 파산 또는 청산할 경우, 채권자들은 프로젝트 회사에 대해 원리금상환을 청구할 수 없다.
  \item 프로젝트 사업주의 도덕적 해이를 방지하기 위해 금융기관은 제한적 소구금융의 장치를 마련해두기도 한다.
\end{choices}
\end{question}
\answer{5}
